\subsection{The MANO Hand Model}\label{subsec:mano}

\acrshort{mano} (hand \textbf{M}odel with \textbf{A}rticulated and \textbf{N}on-rigid
def\textbf{O}rmations) is a statistical, learned model of the human hand
introduced by Romero et al. \cite{MANO:SIGGRAPHASIA:2017}.
It is trained from around 1{,}000 high-resolution 3D scans of 31 subjects
captured across a wide variety of hand poses, totalling 2{,}018 scans of
which 1{,}554 registrations were deemed successful.
Like the SMPL body model \cite{SMPL:2015} on which it is built, \acrshort{mano}
factors geometric changes into those due to the identity of the subject and
those caused by pose.
The resulting model is realistic, low-dimensional, compatible with standard
graphics packages, and capable of representing any human hand.

\subsubsection{Model Formulation}

\acrshort{mano} is based on the SMPL formulation, presented here for completeness
following \cite{SMPL:2015}.
A Linear Blend Skinning (LBS) function $W$ is applied to an articulated
rigged mesh:
\begin{equation}\label{eq:mano_full}
  M(\vec{\boldsymbol{\beta}},\,\vec{\boldsymbol{\theta}})
  = W\!\left(
      T_P(\vec{\boldsymbol{\beta}},\vec{\boldsymbol{\theta}}),\;
      J(\vec{\boldsymbol{\beta}}),\;
      \vec{\boldsymbol{\theta}},\;
      \mathcal{W}
    \right),
\end{equation}
where $\vec{\boldsymbol{\beta}}$ are the shape parameters,
$\vec{\boldsymbol{\theta}}$ are the pose parameters,
$J(\vec{\boldsymbol{\beta}})$ are the joint locations defining the kinematic
tree, and $\mathcal{W}$ are the blend weights.
The posed mesh shape $T_P$ is defined as
\begin{equation}\label{eq:mano_tp}
  T_P(\vec{\boldsymbol{\beta}},\,\vec{\boldsymbol{\theta}})
  = \bar{T} + B_S(\vec{\boldsymbol{\beta}}) + B_P(\vec{\boldsymbol{\theta}}),
\end{equation}
where $\bar{T}$ is the mean template shape, $B_S(\vec{\boldsymbol{\beta}})$
are the shape blend shapes encoding identity variation, and
$B_P(\vec{\boldsymbol{\theta}})$ are the pose blend shapes capturing
pose-dependent mesh deformations.

The pose and shape blend shapes are each defined as a linear combination
of vertex offsets:
\begin{align}
  B_P(\vec{\boldsymbol{\theta}};\,\mathcal{P})
  &= \sum_{n=1}^{9K}
     \Bigl(R_n(\vec{\boldsymbol{\theta}}) - R_n(\vec{\boldsymbol{\theta}}^*)\Bigr)
     \mathbf{P}_n,
  \label{eq:mano_bp}\\[4pt]
  B_S(\vec{\boldsymbol{\beta}};\,\mathcal{S})
  &= \sum_{n=1}^{|\vec{\boldsymbol{\beta}}|}
     \beta_n\,\mathbf{S}_n,
  \label{eq:mano_bs}
\end{align}
where $\mathbf{P}_n \in \mathcal{P}$ are the pose blend shapes, $K$ is the
number of parts in the hand model, $R_n(\vec{\boldsymbol{\theta}})$ indexes
the $n$\textsuperscript{th} element of the concatenated rotation matrix
elements, $\vec{\boldsymbol{\theta}}^*$ is the zero pose, $\beta_n$ are
linear shape coefficients, and $\mathbf{S}_n \in \mathcal{S}$ are the
principal components of the learned shape basis.
The joint locations $J(\vec{\boldsymbol{\beta}})$ also depend on shape and
are learned as a sparse linear regression from mesh vertices, as in SMPL.

Traditional LBS models are overly smooth and suffer from collapse at the
joints. \acrshort{mano} addresses this by learning corrective pose blend shapes that
correct these artifacts, resulting in more natural-looking finger bending.
